% Part: first-order-logic
% Chapter: proof-systems
% Section: natural-deduction

\documentclass[../../../include/open-logic-section]{subfiles}

\begin{document}

\iftag{FOL}
      {\olfileid{fol}{prf}{ntd}}
      {\olfileid{pl}{prf}{ntd}}

\olsection{স্বাভাবিক নিষ্পাদন}

স্বাভাবিক নিষ্পাদন এমন একটি !!a{derivation} পদ্ধতি, যার উদ্দেশ্য বাস্তব
যুক্তিপ্রয়োগের প্রতিরূপ দেওয়া (বিশেষত গণিতবিদেরা যে ধরনের নিয়মবদ্ধ
যুক্তিপ্রয়োগ করেন)। বাস্তব যুক্তিপ্রয়োগ কয়েকটি “স্বাভাবিক” ছক অনুসারে
এগোয়। যেমন, ক্ষেত্রবিচারে প্রমাণে একটি বিয়োজনমূলক পূর্বধারণা থেকে সিদ্ধান্ত
প্রতিষ্ঠা করা যায়—বিয়োজনটির যেকোনো পদ থেকেই সিদ্ধান্তটি অনুসৃত হয় দেখিয়ে।
পরোক্ষ প্রমাণে কোনো সিদ্ধান্তের নঞর্থক রূপ স্ববিরোধে পৌঁছায় দেখিয়ে সিদ্ধান্তটি
প্রতিষ্ঠা করা যায়। শর্তাধীন প্রমাণে পূর্বপদ থেকে উত্তরপদ অনুসৃত হয় দেখিয়ে
“যদি \dots তবে \dots” আকারের শর্তাধীন দাবি প্রতিষ্ঠা করা হয়। স্বাভাবিক
নিষ্পাদন এই স্বাভাবিক অনুমানগুলির কয়েকটির আনুষ্ঠানিক রূপ। প্রতিটি যৌক্তিক
সংযোজক ও পরিমাণসূচকের দুটি করে বিধি থাকে—একটি প্রবর্তন-বিধি ও একটি
অপসারণ-বিধি—এবং প্রত্যেকটি এই ধরনের কোনো স্বাভাবিক অনুমান-ছকের সঙ্গে মেলে।
যেমন, $\Intro{\lif}$ শর্তাধীন প্রমাণের এবং $\Elim{\lor}$ ক্ষেত্রবিচারে
প্রমাণের সঙ্গে মেলে। বিশেষভাবে সরল একটি বিধি হলো $\Elim{\land}$; এটি
$!A \land !B$ থেকে~$!A$ (অথবা $!B$) অনুমান করতে দেয়।

অন্য !!{derivation} পদ্ধতি থেকে স্বাভাবিক নিষ্পাদনকে পৃথক করে এমন একটি
বৈশিষ্ট্য হলো অনুমিতির ব্যবহার। স্বাভাবিক নিষ্পাদনে কোনো
!!^a{derivation} হলো !!{formula}s-এর একটি বৃক্ষ। !!{formula}s-এর
বৃক্ষের মূলে একটি !!{formula} থাকে, আর বৃক্ষের “পাতা”গুলি হলো সেই
!!{formula}s, যাদের থেকে সিদ্ধান্তটি নিষ্পাদিত হয়। স্বাভাবিক নিষ্পাদনে
কিছু পাতার !!{formula}s,~!!{derivation}-এর মধ্যে ভূমিকা পালন করলেও
!!{derivation} সিদ্ধান্তে পৌঁছানোর আগেই “ব্যবহৃত হয়ে যায়”। এটি বাস্তব
যুক্তিপ্রয়োগে এমন পূর্বধারণা প্রবর্তনের রীতির সঙ্গে মেলে, যেগুলি শুধু অল্প
সময়ের জন্য কার্যকর থাকে। যেমন, ক্ষেত্রবিচারে প্রমাণে আমরা বিয়োজনটির
প্রতিটি পদের সত্যতা ধরে নিই; শর্তাধীন প্রমাণে পূর্বপদের সত্যতা ধরে নিই;
পরোক্ষ প্রমাণে সিদ্ধান্তটির নঞর্থক রূপের সত্যতা ধরে নিই। কোনো মধ্যবর্তী
ধাপ প্রতিষ্ঠার জন্য এইভাবে সাময়িক অনুমিতি প্রবর্তন করে পরে সেগুলিকে
সরিয়ে দেওয়া স্বাভাবিক নিষ্পাদনের বিশিষ্ট লক্ষণ। স্বাভাবিক নিষ্পাদনের
!!{derivation}-এর পাতায় থাকা সূত্রগুলিকে অনুমিতি বলা হয়, এবং কিছু
অনুমান-বিধি সেগুলিকে “!!{discharge}” করতে পারে। যেমন,~$!A$-সহ কয়েকটি
অনুমিতি থেকে যদি আমাদের~$!B$-এর একটি !!a{derivation} থাকে, তবে
$\Intro{\lif}$ বিধি~$!A \lif !B$ অনুমান করতে এবং~$!A$ রূপের যেকোনো
অনুমিতিকে অবমুক্ত করতে দেয়। (কোন অনুমান-প্রয়োগে কোন অনুমিতি অবমুক্ত হয় তার
হিসাব রাখতে, অনুমানটি এবং তাতে অবমুক্ত অনুমিতিগুলিকে একটি সংখ্যা দিয়ে
চিহ্নিত করি।) !!{derivation}-এর শেষে যে অনুমিতিগুলি !!{undischarged}
থাকে, সেগুলি একত্রে সিদ্ধান্তটির সত্যতার জন্য যথেষ্ট; তাই কোনো
!!a{derivation} প্রতিষ্ঠা করে যে তার !!{undischarged} অনুমিতিগুলি থেকে
তার সিদ্ধান্তটি অর্থগতভাবে অনুসৃত হয়।

স্বাভাবিক নিষ্পাদনভিত্তিক $\Gamma \Proves !A$ সম্পর্কটি তখনই সত্য, যখন
এমন একটি !!a{derivation} আছে যাতে বৃক্ষের শেষ !!{sentence} হলো~$!A$,
এবং প্রতিটি !!{undischarged} পাতা~$\Gamma$-র অন্তর্গত। স্বাভাবিক
নিষ্পাদনে $!A$ তখনই একটি উপপাদ্য, যখন এমন একটি !!a{derivation} আছে যাতে
$!A$ শেষ !!{sentence} এবং সব অনুমিতি !!{discharged}। যেমন, নিচের
!!a{derivation}-টি দেখায় যে $\Proves (!A
\land !B) \lif !A$:
\begin{prooftree}
  \AxiomC{$\Discharge{!A \land !B}{1}$}
  \RightLabel{\Elim{\land}}
  \UnaryInfC{$!A$}
  \DischargeRule{\Intro{\lif}}{1}
  \UnaryInfC{$(!A \land !B) \lif !A$}
\end{prooftree}
লেবেল~$1$ নির্দেশ করে যে \Intro{\lif} অনুমানে $!A \land !B$ অনুমিতিটি
!!{discharged} হয়েছে।

স্বাভাবিক নিষ্পাদনে কোনো সেট~$\Gamma$ তখনই অসঙ্গত, যখন
$\Gamma \Proves \lfalse$। \FalseInt{} বিধির ফলে কোনো অসঙ্গত সেট থেকে
যেকোনো !!{sentence} !!{derive}d করা যায়।

গেরহার্ড গেন্‌ৎসেন ও স্তানিস্‌ওয়াফ ইয়াশকোভ্‌স্কি ১৯৩০-এর দশকে স্বাভাবিক
নিষ্পাদন পদ্ধতি তৈরি করেন; পরে ডাগ প্রাভিৎস ও ফ্রেডেরিক ফিচ পদ্ধতিটির আরও
বিকাশ ঘটান। এর অনুমানগুলি প্রমাণের স্বাভাবিক পদ্ধতির প্রতিরূপ দেয় বলে
দার্শনিকেরা এটি পছন্দ করেন। ফিচ-নির্মিত রূপগুলি প্রায়ই যুক্তিবিদ্যার
প্রাথমিক পাঠ্যপুস্তকে ব্যবহৃত হয়। যুক্তিবিদ্যার দর্শনে কখনো কখনো স্বাভাবিক
নিষ্পাদনের বিধিগুলিকেই যৌক্তিক কারকগুলির অর্থ দিতে সক্ষম বলে ধরা হয়েছে
(“প্রমাণতাত্ত্বিক অর্থতত্ত্ব”)।

\end{document}
